\chapter{Avalia\c{c}\~ao}
\label{ch:avaliacao}

\section{Configura\c{c}\~ao Experimental}

Os experimentos foram executados em um cluster de cinco n\'os. A
Tabela~\ref{tab:resultados} resume a vaz\~ao sob taxas crescentes de
parti\c{c}\~ao.

\begin{table}[htbp]
  \centering
  \caption{Vaz\~ao (requisi\c{c}\~oes/s) sob taxas simuladas de parti\c{c}\~ao.}
  \label{tab:resultados}
  \begin{tabular}{@{}lrr@{}}
    \toprule
    Taxa de parti\c{c}\~ao & Linha de base & Proposto \\
    \midrule
    0\%  & 12.400 & 12.100 \\
    5\%  & 9.800  & 11.600 \\
    10\% & 6.200  & 10.900 \\
    \bottomrule
  \end{tabular}
  \fonte{Elaborado pelo autor.}
\end{table}

\section{Lat\^encia sob Carga}

A Figura~\ref{fig:latencia} mostra a lat\^encia mediana das
requisi\c{c}\~oes conforme a taxa de parti\c{c}\~ao aumenta.

\begin{figure}[htbp]
  \centering
  \rule{0.6\textwidth}{4cm}
  \caption{Lat\^encia mediana em fun\c{c}\~ao da taxa de parti\c{c}\~ao simulada.}
  \label{fig:latencia}
  \fonte{Elaborado pelo autor.}
\end{figure}

\section{Discuss\~ao}

O protocolo proposto mant\'em uma vaz\~ao maior que a linha de base
assim que a taxa de parti\c{c}\~ao ultrapassa 5\%, ao custo de um
pequeno aumento constante de lat\^encia \cite{ongaro2014search}.
